\section{Bridges, value ceiling and safety rails}
\label{sec:bridges}

\subsection{No bridge on day one}

Konstellation launches \textbf{without a bridge} (\dref{D8}). No canonical
token bridge to Ethereum or any other network exists at genesis, and no IBC
channel carries value until governance has set a rate limit on it
(\S\ref{sec:ratelimit}). The consequence is that the total value on the chain at
launch is bounded by what is allocated at genesis --- a \textbf{value ceiling}
that limits the blast radius of any defect while the chain soaks under real
conditions (launch-sequence phase~9, \S\ref{sec:launch}).

Bridges are where Cosmos EVM chains have lost the most money. A bridge contract
is the highest-priority audit item on this chain's list precisely because,
once one exists, it is the largest pool of value reachable from a single bug
(ENGINEERING \S12). Shipping it later, audited, capped and after a soak is a
deliberate trade of early liquidity for a smaller worst case.

\subsection{The sequence for opening one}
\label{sec:bridgeseq}

The steps once the chain is live, in order (\dref{D8}):

\begin{enumerate}
\item \textbf{Soak.} \code{konstellation-1} runs under the value ceiling for a
  month (\S\ref{sec:launch}, phase~9).
\item \textbf{Build and audit, in parallel with the soak.} The bridge contract
  and the IBC rate-limiting middleware are built and audited. (The rate limiter
  is already built; see \S\ref{sec:ratelimit}. The bridge is not.)
\item \textbf{Calibrate caps.} Hard daily and total caps are set using the
  usage signal from the soak period.
\item \textbf{Open, with caps enforced in the contract.} The bridge opens only
  once the soak is clean, both components are audited, and the caps are set.
  Caps are enforced in the contract, never in a frontend.
\end{enumerate}

\todo{Bridge design is undecided: the counterpart network(s), the trust model
(light client, multisig, third-party messaging layer), the operator, and the
initial cap values. \dref{D8} records only that none of it ships on day one and
the order in which it will.}

\subsection{Safety rails}
\label{sec:rails}

Four controls ship before any user money moves (ENGINEERING \S13). They are
independent of bridge timing and are in place at genesis.

\subsubsection{Circuit breaker}
\label{sec:circuit}

The Cosmos SDK circuit breaker (\code{contrib/x/circuit}, \dref{D14}) can
disable specific message types without halting the chain. It is checked at the
message router, in the ante handler and in the mempool pre-check, so a disabled
message is refused at submission with the reason; the check walks into
\code{authz}-nested messages, so a disabled type cannot be smuggled into a
block inside a \code{MsgExec}; and it wraps the transaction-path precompiles,
which call keepers directly and never see the router, so that disabling an
IBC transfer also closes the ICS-20 precompile to every contract and EOA.

Two properties are worth stating for the reader assessing trust:

\begin{itemize}
\item \textbf{The breaker cannot weld itself shut.} It never disables its own
  messages or governance's, whatever the disable list says, and a trip naming
  one of them is refused. A mistaken or compromised administrator cannot lock
  out the reset.
\item \textbf{Authority is governance.} A 3-of-5 operations multisig holds
  super-admin permission, granted as a genesis entry, so the fastest response to
  an advisory is a multisig transaction rather than a three-day vote; governance
  can revoke that permission. \todo{The mainnet operations multisig signer set
  is not yet chosen; it is a genesis entry, not code.}
\end{itemize}

In operational terms: disabling \code{MsgTransfer} stops all IBC sends, Cosmos
and EVM; disabling \code{MsgEthereumTx} pauses the whole EVM. The breaker is
also the mechanism of validator admission: \code{MsgCreateValidator} ships
disabled in genesis state (\dref{D16}, \S\ref{sec:admission}).

The Cosmos SDK deprecated this package in v0.54 and no longer maintains it.
Konstellation uses it anyway --- it is about 1\,300 stable lines and the
application hook it needs is first-class --- and the build fails with a
re-check checklist on every SDK bump; if a later SDK drops the package it will
be vendored into the chain and enter the audit scope (ENGINEERING \S7.1).

\subsubsection{IBC rate limiting}
\label{sec:ratelimit}

The highest-value single control on the list: it converts ``drained'' into
``lost one window''. No rate-limiting module exists for \code{ibc-go} v11, so
Konstellation ships its own \code{x/ratelimit}, roughly 800 lines modelled on
the semantics of the Cosmos community's earlier module (\dref{D15}). It sits
outermost on the transfer stack, for IBC v1 and v2 alike.

\begin{itemize}
\item Quotas are per (channel, denomination), expressed as a \textbf{percentage
  of the denomination's supply} snapshotted at the start of each window,
  \textbf{optionally capped by an absolute amount}; the lower of the two
  applies. The absolute cap exists because a percentage alone is a weak ceiling
  for the native token --- 1\,\% of \kash{} supply is 1\,\% of the whole chain
  --- and because a foreign token's voucher has no supply before its first
  packet, so only an absolute receive cap makes such a path limitable in
  advance.
\item Quotas are on \textbf{net} flow: a round trip does not consume them, and a
  send that fails or times out within the window is undone.
\item Quotas are \textbf{set and removed only by governance}, per channel, and
  \textbf{no channel carries value on mainnet until governance has set one}
  (launch-sequence phase~9 gate). \code{testnet-1} opens no external channels.
\item The module fails closed on a store error rather than open.
\end{itemize}

Alongside it, an ICS-20 receive gate in the compliance module refuses incoming
transfers to a frozen address; the packet is error-acknowledged and the sending
chain refunds (\S\ref{sec:compliance}).

\subsubsection{Bridge caps}

Enforced in the bridge contract, not the frontend (\S\ref{sec:bridgeseq}). Not
applicable until a bridge exists.

\subsubsection{A rehearsed halt drill}

Executed on \code{testnet-1} before mainnet: a simulated advisory at an awkward
hour, a halt at a height, a patch, a coordinated restart, back in under 90
minutes. Repeated quarterly thereafter and written up each time (ENGINEERING
\S13, \S17). The first mainnet incident must not be the first rehearsal.
